\section{Experimental Evaluation}
\label{sec:evaluation}

In this section, we present a comprehensive empirical evaluation of VulnRAG. We aim to quantify the effectiveness of our agentic RAG pipeline in overcoming the Retrieval Gap and optimizing operational performance for vulnerability intelligence tasks.

\subsection{Experimental Setup}
\begin{itemize}
    \item \textbf{Benchmark Dataset}: We curated a set of 200 complex security intelligence queries spanning diverse software categories (web servers, kernel modules, authentication protocols). Each query is mapped to a ground-truth set of operationally linked CVEs.
    \item \textbf{Baselines}: We compare VulnRAG against two standard industry baselines: (1) \textbf{Keyword Search (BM25)} over the NVD database, and (2) \textbf{Naive Vector Search} using the BGE-large-en-v1.5 embedding model.
    \item \textbf{Models}: Our tiered routing evaluation utilizes Gemini 1.5 Flash (Flash), Gemini 1.5 Pro (Pro), and Claude 3.5 Sonnet (Reasoning) as representative LLM tiers.
\end{itemize}

\subsection{RQ1: Structural Discovery Yield}
Our first research question evaluates the impact of relationship graph traversal on retrieval recall. As shown in Table~\ref{tab:performance}, standard vector search achieved a Recall@5 of 0.684. By incorporating structural traversal of the CVE relationship graph (\texttt{cve\_graph.py}), VulnRAG increased recall to \textbf{0.892}. Furthermore, Fig.~\ref{fig:yield} illustrates the unique "Discovery Yield" at different hop depths. Traversal at depth 2 revealed 2.1$\times$ more relevant vulnerabilities than naive semantic search, particularly for vulnerabilities linked via shared libraries or vendor-specific patches that lacked textual overlap.

\subsection{RQ2: Impact of Agentic Self-Correction}
RQ2 assesses the effectiveness of the zero-shot Evaluator-Reformulator feedback loop. As detailed in Table~\ref{tab:correction}, the initial retrieval pass for complex, ambiguous queries often suffered from low precision (0.592). However, upon triggering the agentic self-correction loop, the Query Reformulator successfully expanded product scopes and extracted missing CPE identifiers. This iterative process pushed the Precision-Recall frontier (Fig.~\ref{fig:pr_frontier}), matching the performance of fine-tuned reflective models without requiring task-specific training data.

\subsection{RQ3: Reranking and Intelligence Density}
To address RQ3, we evaluated the Multi-Factor Semantic Reranking (MFSR) algorithm. By integrating CVSS severity and exploit availability into the ranking logic (Table~\ref{tab:strategies}), VulnRAG achieved an NDCG@5 score of \textbf{0.854}, a significant improvement over the naive vector baseline (0.642). This improvement indicates that MFSR effectively prioritizes "Actionable Intelligence," ensuring that the most critical and exploitable vulnerabilities are placed in the high-attention regions of the LLM context, thereby mitigating the "Lost in the Middle" problem.

\subsection{RQ4: Operational Economics and Tiered Routing}
Finally, we analyzed the performance-to-cost trade-off of our complexity-aware multi-tier routing. Fig.~\ref{fig:economics} plots reasoning success rate against operational inference cost. While the All-Flagship configuration achieved the highest quality, it incurred significant financial overhead. VulnRAG’s routing engine successfully escalated only the most complex 18\% of tasks to the Reasoning tier, while handling routine lookups via the Flash tier. This strategy achieved a 97.4\% success rate at an \textbf{84.2\% lower cost} than the flagship-only approach, demonstrating the economic viability of agentic VAPT at enterprise scale.
